\documentclass[12pt]{article}

\usepackage[a4paper, total={6in, 8in}]{geometry}
\usepackage{textcomp}

\begin{document}
\setlength{\parindent}{0pt}

\def \t {\textrightarrow}
\def \v {\vspace{1em}}
\def \bi {\begin{itemize}}
\def \ei {\end{itemize}}
\def \s[#1] {\section*{#1}}
\def \ss[#1] {\subsection*{#1}}
\def \sss[#1] {\subsubsection*{#1}}

\s[Carducci - Alla stazione]
Stazione di Bologna \t che permette di individuare un focus sull'attualità con la strage di Bologna del 1980 \\
Clima del distacco, di insofferenza \t atmosfera simile a quella dello spleen \\
La donna deve fare ritorno da suo marito \t compresenza di due figure, moglie e amante \t la donna viene vissuto nel mondo del presente e nel mondo dei ricordi, che si colorano \t capelli castanei \\
Un giugno che scalda ma non è ancora infuocato \t tepore di giugno \\
È il canto del distacco, che pone al centro la figura femminile, che accompagna la duplicità di questa unica figura femminile \t e il ritorno allo stato di origine \\
Il progresso allontana \\
Siamo nel 73 (o 63) in cui Carolina Cristofori Piva (che viene chiamata Lidia qua) \t il nome Lidia è presente nella cultura classica

\v

Inizia subito con esclamazione \\
I fanali sono i lampioni che si ripetono ripetitivi dietro agli alberi, i cui rami sono intrisi di pioggia (evoca pioggia del pineto d'annunzio) \\
Luce che sbadiglia, sinestesia \t stridual è onomatopea \\
Flebile acuta etc. climax \t finisce con fischia \t il mondo del progresso è la vaporiera, che rappresenta il progresso del tempo \\
Chi resta e chi va, disallineamento \t BOCCIONI \\
Cielo plumbeo ripreso direttamente da Boudelaire \t che soffoca, un tedio che caratterizza l'universo dei maledetti \\
Analogia, come un grande fantasma \\
Verso cosa e dove tende questa gente, che si affretta verso i carri scuri (le carrozze del teno) \t una coralità, che genera un unico personaggio \t tutti sono uno, che sono accomunati da qualcosa che cercano \\
Ravvolta = coperti, per il freddo \\
Incomunicabilità presente, incognita \\
Tu e ora \t va dall'esterno all interno del suo animo \\
Istanti gioiti, ... climax con apice ricordi (focus introspettivo), invece tempo incalzante (focus sulla sfera esterna)

\v

Che spasmo pare = anastrofe \\
Incappucciati di nero \t sono le guardie \t analogia, come ombre \\
Verso 19 si interrompe a meta con un segno di interp. semi forte, per continuare la descrizione di un universo spettrale \t quasi come in un racconti di Alan Poe \\
Sono gli attrezzi per rincardinare le rotaie \t la presenza del ferreo (queste mazze, ferrei freni, macchina di ferro) \t ferro e ferrei è una figura etimolgoica, no poliptoto \\
Lugubre rintocco evoca presagio di morte ed evoca Boudelaire, in cui l'anima sente e risponde, ha dentro di se uno scuotimento dettato da questo tedio \\
Il tedio carducciano è boudelriano, questo non senso \t che è diverso dal tedio leopardiano

\v

Dal 25 al 30 \t gli sportelli chiusi fanno violenza all'animo che rimane chiuso \\
Scroscia è onomatopea \\
L'empio mostro è l'elemento identificativo di questa vaporiera \t empio, malvagio \t l'emblema del progresso è un mostro \t perchè allontana \\
Divorati dal vuoto \\
Nel verso 34 c'è espressione portasi \t si porta via lidia, la speranza, l'enigma di un non ben identificato ritorno \\
Bianca faccia pallido velo ? \t chiasmo \t il velo è un fazzoletto, ed era consuetidine salutare con il fazzoletto \\
Dal verso 37 alla fine si ha una porzione interamente dedicata alla donna \t la bianca faccia è quella di Lidia, e il focus è su di lei fino al verso 53

\v

L'architettura del testo è ben individuabile \t il testo unisce l'esterno e l'interno, il paesaggio lugubre e lo stato d'animo \\
Il vuoto che inghiottisce il treno è lo stesso che inghiottisce l'animo del poeta \\
Lettura verghiana con l'ostrica attaccata allo scoglio \t l'allontanarsi annulla la nostra identità, destruttura \t qua sofferenza interiorizzata
\end{document}
